\section{模型检验}

\subsection{检验体系总览}

\begin{table}[H]
\centering\small
\caption{检验清单：方法关系、判据与结果}
\begin{tabular}{p{3.6cm}p{2.1cm}p{3.4cm}p{4.6cm}}
\toprule
检验对 & 关系类型 & 判据 & 结果 \\
\midrule
CRITIC vs PCA/TOPSIS/等权 & 互相验证 & 样本级 Spearman & $0.871/0.946/0.991$，排序稳健 \\
clr 线性 vs 梯度提升 & 互相验证 & $10^5$ 采样点秩相关 & $0.77$，$\bm p^*$ 非模型形式幻觉 \\
clr vs 裸 OLS & 改进证据 & 检验集 $R^2$/RMSE & $0.772$ vs $0.654$，闭合效应实证 \\
A1 vs A2/A3 全量 & 代表性检验 & KS、Cohen's $d$ & $p=0.60/0.12$，$d\approx0.005$ \\
冲突结论跨集复现 & 稳健性 & 冲突率与主类型 & $21.5\%/21.1\%/22.9\%$，类型一致 \\
1M$\to$60M$\to$1B & 假设~\ref{asm:proxy} 检验 & 排名 Spearman & $0.52/0.47/0.68$；截距漂、斜率稳 \\
10B/70B 外推表 & 稳健性讨论 & 修正前后 MAE & $3.5\to0.10$（非观测，只作讨论） \\
线性乘性耦合 vs 有效数据 & 形式选择 & 五折 cv-RMSE & $0.0497$ vs $0.0979$ \\
解析解 vs 数值优化 & 互相验证 & 相对误差 & $<10^{-6}$（命题 2、4） \\
标度律 vs 真实文献点 & 外部合理性 & 趋势 Spearman & $0.959$--$0.983$，仅比趋势 \\
质量分 vs 训练价值 & 概念区分 & $\rho_{Q\nu}$ & $-0.06$：分歧本身即结论 \\
\bottomrule
\end{tabular}
\label{tab:checks}
\end{table}

三点说明。其一，"互相验证"仅用于\emph{同一问题、不同假设}的方法对（如 CRITIC 与 PCA），判据取排序一致性；分工衔接的上下游方法（如冲突分析与配比回归）之间不存在验证关系，不作此类声称。其二，不一致处同样是结论：$\rho_{Q\nu}=-0.06$ 是"内在质量 $\ne$ 训练价值"的定量证据而非检验失败；60M/1B 的绝对预测失效（$R^2$ 大负值）源于截距漂移，恰好论证了"跨尺度外推必须做尺度修正"并给出修正后 35 倍的误差收窄。其三，A12--A15 与 B8 的证据等级（外推生成/不同源）全程继承标注，未升格为实证。

\subsection{灵敏度分析}

\textbf{冲突阈值}：$\tau\in[0.6,0.85]$ 内冲突率 $12\%$--$35\%$ 连续变化，域排序与主冲突类型不变。\textbf{评价权重}：等权与 pile\_cc 单域两口径下 $\bm p^*$ 的头部域一致（pile\_cc 居首），仅集中度不同。\textbf{质量口径锚点}（假设~\ref{asm:anchor}）：主口径 $Q_0=0.6609$ 与副口径 $Q_{\mathrm B}=\bar q/0.6609$ 下，$Q_{\mathrm B}:0.5\to0.6$ 的等效倍数为 $1.29/1.42/1.77$ 与 $1.47/1.73/2.48$（$C=10^{20}/10^{22}/10^{24}$），方向不变。\textbf{消解规则}：Huber 降权与调和平均两口径下域级质量排序 Spearman $>0.95$。\textbf{bootstrap}：$\kappa_N,\kappa_D,k_0$ 的 95\% CI 为 $[0.390,0.473]$、$[0.085,0.195]$、$[0.114,0.168]$，$P(\kappa_N>\kappa_D)=1$。

\section{模型的优缺点与改进方向}

\textbf{优点}：(i) 全流程证据分层，半合成与外推数据未混入主结论；(ii) 冲突率、闭合效应、质量—价值解耦三处以量化证据修正了《数据说明》页眉与文献的现成说法；(iii) 广义标度律的形式由交叉验证从五种候选中裁决，且四条解析命题全部闭式可验；(iv) 问题一至问题二的接口（$Q$ 口径、$\bm p^*$、迁移矩阵、漂移律）落盘为文件，问题三、四可直接消费。

\textbf{缺点与改进}：(i) 17 域质量分中 11 个 inferred 域依赖校准函数（$R^2_{\mathrm{cv}}=0.31$），信息量有限，若增补这些域的模型评分器信号可显著收紧区间；(ii) 质量交互项 $\lambda$ 符号域间不一致的机制尚未解释，需按验证域类型分层建模；(iii) B 附件为半合成数据，$\kappa_N,\kappa_D,k_0$ 的绝对数值不宜直接外推到真实训练体系；(iv) $s(N)$ 只在 1M--1B 上标定，1B 以上是外推。

\section{结论}

问题一建立了从 22 个异质信号到连续质量分的可复现管线，测得并消解了 $21.5\%$ 的指标冲突，以 clr 回归给出 17 域配比的迁移矩阵与最优配比（预测降损 $9.9\%$），并发现内在质量与训练价值近乎正交。问题二以数据裁决的线性乘性耦合建立广义标度律，主口径 $Q_0=0.6609$。四条解析推论——质量弹性约为参数弹性的 $3.7$ 倍、等效替代随规模从 $1.29$ 倍放大到 $1.77$ 倍、质量地板的存在、最优配置随质量向"小模型多数据"移动——共同刻画了"数据质量是与算力并列的一等资源"这一核心判断。

\begin{table}[H]
\centering\small
\caption{问题一、二产出与下游接口}
\begin{tabular}{lll}
\toprule
产出 & 内容 & 下游用途 \\
\midrule
P1-Q\_sample / Q\_domain & 27.2 万样本级 $Q$；17 域 $\hat Q_d$ 及 CI & 问题二 $\bar Q(\bm p)$；问题三 $Q_0$ \\
P1-$f(\bm p)$ / $T$ & 配比回归系数；$17\times13$ 迁移矩阵 & 问题三损失项；领域互补分析 \\
P1-$\bm p^*$ & 双口径最优配比，$\bar Q(\bm p^*)$ & 问题三默认配比 \\
P1-rank\_decay & 排名保持度随 $N$ 的三点曲线 & 外推不确定性口径 \\
P2-params / theory & 式~\eqref{eq:main} 全参数与 CI；$N^*(C,Q)$ 闭式 & 问题三目标函数与优化初值 \\
\bottomrule
\end{tabular}
\label{tab:iface}
\end{table}

\begin{thebibliography}{9}\small
\bibitem{hoffmann} Hoffmann J, et al. Training Compute-Optimal Large Language Models. NeurIPS 2022.
\bibitem{regmix} Liu Q, et al. RegMix: Data Mixture as Regression for Language Model Pre-training. ICLR 2025.
\bibitem{metarater} SlimPajama-Meta-rater 数据集（OpenDataLab）：多维质量信号体系.
\bibitem{fineweb} Penedo G, et al. The FineWeb Datasets. NeurIPS 2024（FineWeb-Edu 教育价值过滤）.
\bibitem{muennighoff} Muennighoff N, et al. Scaling Data-Constrained Language Models. NeurIPS 2023.
\bibitem{aitchison} Aitchison J. The Statistical Analysis of Compositional Data. Chapman \& Hall, 1986.
\bibitem{critic} Diakoulaki D, Mavrotas G, Papayannakis L. Determining objective weights in multiple criteria problems: The CRITIC method. Computers \& Operations Research, 1995.
\bibitem{huber} Huber P J. Robust estimation of a location parameter. Ann. Math. Statist., 1964.
\bibitem{pythia} Biderman S, et al. Pythia: A Suite for Analyzing Large Language Models. ICML 2023.
\end{thebibliography}

\appendix
\section*{附录：AI 使用披露}
依据赛题规定，AI 工具仅用于资料整理、代码调试与文字排版辅助；本文的模型选择、公式推导、结果分析与论证由参赛队完成并逐项复核。全部数值结果可由随文代码（\texttt{code\_q1/}、\texttt{code\_q2/}）在附件数据上复现。
